\documentclass[10pt,a4paper]{amsart}
\usepackage{amsmath,amssymb,amsthm,graphicx}
% Retained preamble for the Stacks Project chapter {CH} (shared project preamble).
% The chapter file pinned-original-source.tex only carries \input{{preamble}} on its first line,
% so no preamble text lives inside the chapter itself.  The Stacks Project uses one single shared
% preamble for every chapter: the snapshot stages it three times (sources/stacks-chow/preamble.tex,
% sources/stacks-etale-cohomology/preamble.tex and focused-seeds/AGP/sources/stacks-snapshot-20261011/
% preamble.tex) and those copies are identical in size.  The block below is copied mechanically,
% line for line, from sources/stacks-chow/preamble.tex lines 14-18 and 154-217: the Xy-pic setup and
% the native theorem declarations together with the mathematical macros.  Nothing was added, removed
% or redefined; the wrapper class and the amsmath/amssymb/amsthm/graphicx base packages are supplied
% by the assembler.
% For commutative diagrams we use Xy-pic
% Chapter: stacks-spaces-more-morphisms.tex -- More on Morphisms of Spaces.
\usepackage[all]{xy}

% We use 2cell for 2-commutative diagrams.
\xyoption{2cell}

% Theorem environments.
%
\theoremstyle{plain}
\newtheorem{theorem}[subsection]{Theorem}
\newtheorem{proposition}[subsection]{Proposition}
\newtheorem{lemma}[subsection]{Lemma}

\theoremstyle{definition}
\newtheorem{definition}[subsection]{Definition}
\newtheorem{example}[subsection]{Example}
\newtheorem{exercise}[subsection]{Exercise}
\newtheorem{situation}[subsection]{Situation}

\theoremstyle{remark}
\newtheorem{remark}[subsection]{Remark}
\newtheorem{remarks}[subsection]{Remarks}

\numberwithin{equation}{subsection}

% Macros
%
\def\lim{\mathop{\mathrm{lim}}\nolimits}
\def\colim{\mathop{\mathrm{colim}}\nolimits}
\def\Spec{\mathop{\mathrm{Spec}}}
\def\Hom{\mathop{\mathrm{Hom}}\nolimits}
\def\Ext{\mathop{\mathrm{Ext}}\nolimits}
\def\SheafHom{\mathop{\mathcal{H}\!\mathit{om}}\nolimits}
\def\SheafExt{\mathop{\mathcal{E}\!\mathit{xt}}\nolimits}
\def\Sch{\mathit{Sch}}
\def\Mor{\mathop{\mathrm{Mor}}\nolimits}
\def\Ob{\mathop{\mathrm{Ob}}\nolimits}
\def\Sh{\mathop{\mathit{Sh}}\nolimits}
\def\NL{\mathop{N\!L}\nolimits}
\def\CH{\mathop{\mathrm{CH}}\nolimits}
\def\proetale{{pro\text{-}\acute{e}tale}}
\def\etale{{\acute{e}tale}}
\def\QCoh{\mathit{QCoh}}
\def\Ker{\mathop{\mathrm{Ker}}}
\def\Im{\mathop{\mathrm{Im}}}
\def\Coker{\mathop{\mathrm{Coker}}}
\def\Coim{\mathop{\mathrm{Coim}}}

% Boxtimes
%
\DeclareMathSymbol{\boxtimes}{\mathbin}{AMSa}{"02}

%
% Macros for moduli stacks/spaces
%
\def\QCohstack{\mathcal{QC}\!\mathit{oh}}
\def\Cohstack{\mathcal{C}\!\mathit{oh}}
\def\Spacesstack{\mathcal{S}\!\mathit{paces}}
\def\Quotfunctor{\mathrm{Quot}}
\def\Hilbfunctor{\mathrm{Hilb}}
\def\Curvesstack{\mathcal{C}\!\mathit{urves}}
\def\Polarizedstack{\mathcal{P}\!\mathit{olarized}}
\def\Complexesstack{\mathcal{C}\!\mathit{omplexes}}
% \Pic is the operator that assigns to X its picard group, usage \Pic(X)
% \Picardstack_{X/B} denotes the Picard stack of X over B
% \Picardfunctor_{X/B} denotes the Picard functor of X over B
\def\Pic{\mathop{\mathrm{Pic}}\nolimits}
\def\Picardstack{\mathcal{P}\!\mathit{ic}}
\def\Picardfunctor{\mathrm{Pic}}
\def\Deformationcategory{\mathcal{D}\!\mathit{ef}}
\usepackage{cleveref}
\crefname{originalreference}{original reference}{original references}
\Crefname{originalreference}{Original reference}{Original references}
\makeatletter
\newlabel{etale-lemma-etale-etale-local}{{04HL}{source}{}{sourcepointer-1}{}}
\newlabel{etale-lemma-etale-etale-local@cref}{{[originalreference][0][]04HL}{[1][0][]source}{}{}{}}
\makeatother
\crefname{originallemma}{original lemma}{original lemmas}
\Crefname{originallemma}{Original lemma}{Original lemmas}
\makeatletter
\newlabel{lemma-etale-splits-off-just-one-quasi-finite-part}{{0ADU}{source}{}{sourcepointer-2}{}}
\newlabel{lemma-etale-splits-off-just-one-quasi-finite-part@cref}{{[originallemma][0][]0ADU}{[1][0][]source}{}{}{}}
\makeatother
\crefname{originalreference}{original reference}{original references}
\Crefname{originalreference}{Original reference}{Original references}
\makeatletter
\newlabel{situation-descend-resolution-property}{{0GUY}{source}{}{sourcepointer-3}{}}
\newlabel{situation-descend-resolution-property@cref}{{[originalreference][0][]0GUY}{[1][0][]source}{}{}{}}
\makeatother
\crefname{originalreference}{original reference}{original references}
\Crefname{originalreference}{Original reference}{Original references}
\makeatletter
\newlabel{spaces-cohomology-lemma-finite-rings}{{0DK4}{source}{}{sourcepointer-4}{}}
\newlabel{spaces-cohomology-lemma-finite-rings@cref}{{[originalreference][0][]0DK4}{[1][0][]source}{}{}{}}
\makeatother
\crefname{originalreference}{original reference}{original references}
\Crefname{originalreference}{Original reference}{Original references}
\makeatletter
\newlabel{spaces-cohomology-lemma-homs-over-open}{{07UM}{source}{}{sourcepointer-5}{}}
\newlabel{spaces-cohomology-lemma-homs-over-open@cref}{{[originalreference][0][]07UM}{[1][0][]source}{}{}{}}
\makeatother
\crefname{originalreference}{original reference}{original references}
\Crefname{originalreference}{Original reference}{Original references}
\makeatletter
\newlabel{spaces-descent-lemma-finite-over-finite-module}{{060Z}{source}{}{sourcepointer-6}{}}
\newlabel{spaces-descent-lemma-finite-over-finite-module@cref}{{[originalreference][0][]060Z}{[1][0][]source}{}{}{}}
\makeatother
\crefname{originalreference}{original reference}{original references}
\Crefname{originalreference}{Original reference}{Original references}
\makeatletter
\newlabel{spaces-morphisms-lemma-affine-equivalence-modules}{{08AI}{source}{}{sourcepointer-7}{}}
\newlabel{spaces-morphisms-lemma-affine-equivalence-modules@cref}{{[originalreference][0][]08AI}{[1][0][]source}{}{}{}}
\makeatother
\title{Coherent sheaves on the Noetherian space are quotients of sums of pushforwards of powers of the ideal}
\author{The Stacks Project Authors}
\date{stacks-project:spaces-more-morphisms.tex:local-snapshot-20261011; GNU FDL 1.2 or later}
\begin{document}
\maketitle
\noindent\textit{Editorial note.} Counted span is the author's own proof environment, source lines 13452-13535, immediately following the statement at 13439-13450.\par\medskip
\noindent\textit{Editorial note.} The displayed chain identifying global sections is retained verbatim; it is the technical heart of the proof.\par\medskip
\noindent\textit{Editorial note.} Five cross-chapter pointers (Etale Cohomology, Coherent, Descent on Spaces, Morphisms of Spaces) are recorded in assembly.external\_references.\par\medskip
\noindent\textit{Editorial note.} Steps the author leaves to the reader are flagged here for the reviewer: source line(s) 13519 carry an "omitted"/"follows easily"/"immediate from the definitions" phrase. Those sentences are retained verbatim; the displayed argument around them stays inside the counted span.\par\medskip
\noindent\textit{Allowed standard prerequisites (editorial).}\begin{itemize}
\item The situation of descending the resolution property. chapter situation situation-descend-resolution-property, cited in the statement. Fixes the finite locally free covering pi : Y -> X and the classifying ideal I.
\item Etale splitting off of a single quasi-finite part. chapter lemma lemma-etale-splits-off-just-one-quasi-finite-part, cited in the proof. Produces the etale neighbourhood and the single-point finite piece used to construct local generators.
\item Duality of finite modules over the covering. Derived Categories of Spaces lemmas on finite modules over open pieces and Spaced cohomology lemmas, cited in the proof. Identifies the module of homomorphisms with a global section module so that generators can be transported.
\item The canonical bundle of the finite locally free covering. the sheaf omega\_\{Y/X\} appearing in the proof (dualising module of the covering). Turns sections of pi\textasciicircum{}*F twisted by omega into sections of F, which is how the displayed sum of pushforwards is produced.
\end{itemize}\medskip
\ifdefined\hypertarget\hypertarget{sourcepointer-1}{}\fi
\noindent\textit{Source reference.} Original reference 04HL: focused-seeds/geometry/S9-spaces-surjection-noetherian-case/ref-tags.txt, source lines 4--4. Navigation pointer only: the counted proof cites this statement by \textbackslash{}ref as a standing input; the deduction the counted proof has to make is carried out inside the counted span, and no part of that deduction is delegated to this pointer.\par\medskip
\ifdefined\hypertarget\hypertarget{sourcepointer-2}{}\fi
\noindent\textit{Source reference.} Original lemma 0ADU: sources/stacks-spaces-more-morphisms/pinned-original-source.tex, source lines 7544--7544. Navigation pointer only: the counted proof cites this statement by \textbackslash{}ref as a standing input; the deduction the counted proof has to make is carried out inside the counted span, and no part of that deduction is delegated to this pointer.\par\medskip
\ifdefined\hypertarget\hypertarget{sourcepointer-3}{}\fi
\noindent\textit{Source reference.} Original reference 0GUY: sources/stacks-spaces-more-morphisms/pinned-original-source.tex, source lines 13416--13416. Navigation pointer only: the counted proof cites this statement by \textbackslash{}ref as a standing input; the deduction the counted proof has to make is carried out inside the counted span, and no part of that deduction is delegated to this pointer.\par\medskip
\ifdefined\hypertarget\hypertarget{sourcepointer-4}{}\fi
\noindent\textit{Source reference.} Original reference 0DK4: focused-seeds/geometry/S9-spaces-surjection-noetherian-case/ref-tags.txt, source lines 5--5. Navigation pointer only: the counted proof cites this statement by \textbackslash{}ref as a standing input; the deduction the counted proof has to make is carried out inside the counted span, and no part of that deduction is delegated to this pointer.\par\medskip
\ifdefined\hypertarget\hypertarget{sourcepointer-5}{}\fi
\noindent\textit{Source reference.} Original reference 07UM: focused-seeds/geometry/S9-spaces-surjection-noetherian-case/ref-tags.txt, source lines 6--6. Navigation pointer only: the counted proof cites this statement by \textbackslash{}ref as a standing input; the deduction the counted proof has to make is carried out inside the counted span, and no part of that deduction is delegated to this pointer.\par\medskip
\ifdefined\hypertarget\hypertarget{sourcepointer-6}{}\fi
\noindent\textit{Source reference.} Original reference 060Z: focused-seeds/geometry/S9-spaces-surjection-noetherian-case/ref-tags.txt, source lines 7--7. Navigation pointer only: the counted proof cites this statement by \textbackslash{}ref as a standing input; the deduction the counted proof has to make is carried out inside the counted span, and no part of that deduction is delegated to this pointer.\par\medskip
\ifdefined\hypertarget\hypertarget{sourcepointer-7}{}\fi
\noindent\textit{Source reference.} Original reference 08AI: focused-seeds/geometry/S9-spaces-surjection-noetherian-case/ref-tags.txt, source lines 8--8. Navigation pointer only: the counted proof cites this statement by \textbackslash{}ref as a standing input; the deduction the counted proof has to make is carried out inside the counted span, and no part of that deduction is delegated to this pointer.\par\medskip
\begin{lemma}
\label{lemma-surjection-in-Noetherian-case}
In Situation \ref{situation-descend-resolution-property}, assume
$X$ is Noetherian. Then for any coherent $\mathcal{O}_X$-module
$\mathcal{F}$ there exist $r \geq 0$,
integers $n_1, \ldots, n_r \geq 0$, and a surjection
$$
\bigoplus\nolimits_{i = 1, \ldots, r} \pi_*(\mathcal{I}^{n_i})
\longrightarrow \mathcal{F}
$$
of $\mathcal{O}_X$-modules.
\end{lemma}

\begin{proof}
Denote $\omega_{Y/X}$ the coherent $\mathcal{O}_Y$-module such that there
is an isomorphism
$$
\pi_*\omega_{Y/X} \cong
\SheafHom_{\mathcal{O}_X}(\pi_*\mathcal{O}_Y, \mathcal{O}_X)
$$
of $\pi_*\mathcal{O}_Y$-modules, see Morphisms of Spaces, Lemma
\ref{spaces-morphisms-lemma-affine-equivalence-modules} and
Descent on Spaces, Lemma \ref{spaces-descent-lemma-finite-over-finite-module}.
The canonical map $\mathcal{O}_X \to \pi_*\mathcal{O}_Y$
produces a canonical map
$$
\text{Tr}_\pi : \pi_*\omega_{Y/X} \longrightarrow \mathcal{O}_X
$$
Since $V$ is Noetherian affine we may choose sections
$$
s_1, \ldots, s_r \in
\Gamma(V, \pi^*\mathcal{F} \otimes_{\mathcal{O}_Y} \omega_{Y/X})
$$
generating the coherent module
$\pi^*\mathcal{F} \otimes_{\mathcal{O}_X} \omega_{Y/X}$
over $V$. By
Cohomology of Spaces, Lemma \ref{spaces-cohomology-lemma-homs-over-open}
we can choose integers $n_i \geq 0$ such that $s_i$
extends to a map
$s_i' : \mathcal{I}^{n_i} \to
\pi^*\mathcal{F} \otimes_{\mathcal{O}_Y} \omega_{Y/X}$.
Pushing to $X$ we obtain maps
$$
\sigma_i :
\pi_*\mathcal{I}^{n_i}
\xrightarrow{\pi_*s'_i}
\pi_*(\pi^*\mathcal{F} \otimes_{\mathcal{O}_Y} \omega_{Y/X}) =
\mathcal{F} \otimes_{\mathcal{O}_X} \pi_*\omega_{Y/X}
\xrightarrow{\text{Tr}_\pi} \mathcal{F}
$$
where the equality sign is
Cohomology of Spaces, Lemma \ref{spaces-cohomology-lemma-finite-rings}.
To finish the proof we will show that the sum of these maps is surjective.

\medskip\noindent
Let $x \in |X|$ be a point of $X$. Let $v \in |V|$ be a point mapping to $x$.
We may choose an \'etale neighbourhood $(U, u) \to (X, x)$ such that
$$
U \times_X Y = W \coprod W'
$$
(disjoint union of algebraic spaces) such that $W \to U$ is an
isomorphism and such that the unique point $w \in W$ lying over
$u$ maps to $v$ in $V \subset Y$. To see this is true use
Lemma \ref{lemma-etale-splits-off-just-one-quasi-finite-part} and
\'Etale Morphisms, Lemma \ref{etale-lemma-etale-etale-local}.
After shrinking $U$ further if necessary we may assume $W$ maps into
$V \subset Y$ by the projection. Since the formation of $\omega_{Y/X}$
commutes with \'etale localization we see that
$$
\pi_*\omega_{Y/X}|_U = (\pi|_W)_*\omega_{W/U} \oplus (\pi|_{W'})_*\omega_{W'/U}
$$
We have $(\pi|_W)_*\omega_{W/U} = \mathcal{O}_U$ and this isomorphism
is given by the trace map $\text{Tr}_\pi|_U$ restricted
to the first summand in the decomposition above. Since $W$ maps
into $V$ we see that $\mathcal{I}^{n_i}|_W = \mathcal{O}_W$.
Hence
$$
\pi_*(\mathcal{I}^{n_i})|_U = \mathcal{O}_U \oplus
(W' \to U)_*(\mathcal{I}^{n_i}|_{W'})
$$
Chasing diagrams the reader sees (details omitted) that $\sigma_i|_U$
on the summand $\mathcal{O}_U$ is the map $\mathcal{O}_U \to \mathcal{F}$
corresponding to the section
$$
s_i|_W  \in
\Gamma(W,\pi^*\mathcal{F} \otimes_{\mathcal{O}_Y} \omega_{Y/X})
=
\Gamma(W, \mathcal{F}|_W \otimes_{\mathcal{O}_W} \omega_{W/U})
=
\Gamma(U, \mathcal{F})
$$
Since the sections $s_i$ generate the module
$\pi^*\mathcal{F} \otimes_{\mathcal{O}_Y} \omega_{Y/X}$
over $V$ and since $W$ maps into $V$ we conclude
that the restriction of $\bigoplus \sigma_i$ to $U$ is surjective.
Since $x$ was an arbitrary point the proof is complete.
\end{proof}

\end{document}
